\chapter{Ciele práce}

\label{kap:ciele}


\section{Varianty v DNA}

V prírode sa často vyskytujú dvojice DNA molekúl, ktoré obsahujú veľmi podobnú postupnosť báz, líšiacu sa iba v malých detailoch. Typickým príkladom sú DNA dvoch rôznych jedincov rovnakého druhu. Tieto malé odlišnosti v DNA voláme \emph{varianty}. Najjednoduchšie druhy variantov sú nasledujúce.

\begin{description}
\item [Jednonukleotidový polymorfizmus (SNP\footnotemark).]\footnotetext{z anglického \emph{\textbf{s}ingle \textbf{n}ucleotide \textbf{p}olymorphism}} Jedna báza z prvej DNA postupnosti sa v druhej postupnosti nahradí inou bázou.

\item [Inzercia.] Do postupnosti je vsunutá skupina báz.

\item [Delécia.] Z postupnosti vypadne súvislá skupina báz.

\end{description}

\begin{figure}[t]
\centerline{\includegraphics[width=0.4\textwidth]{images/varianty}}
\caption[Varianty v DNA]{Jednonukleotidový polymorfyzmus, inzercia a delécia.}
\label{obr:varianty}
\end{figure}

\section{Identifikácia variantov}

Pri niektorých využitiach DNA sekvenovania sa sekvenuje vzorka, o ktorej je známe, že by
sa mala podobať na inú, už osekvenovanú DNA. 
Cieľom sekvenovania je potom zistiť, ako sa tieto dve DNA postupnosti líšia. Jedným z takýchto využití 
je napríklad zisťovanie rezistencie baktérií na antibiotiká \cite{Bradley2015}.

V našej práci sa budeme zaoberať nasledujúcim scenárom. Máme nejakú známu postupnosť
dusíkatých báz, ktorú budeme nazývať \emph{referencia}. Ďalej máme vzorku DNA, o ktorej vieme,
že sa od referencie líši len veľmi málo. Táto vzorka bola spracovaná prístrojom MinION, máme teda k dispozícii
nameraný surový signál z jednotlivých čítaní. Naším cieľom je identifikovať varianty v sekvenovanej
vzorke vzhľadom na referenciu. Ideálne by bolo vedieť s dobrou presnosťou určovať varianty už z jedného čítania.

Snažíme sa teda navrhnúť algoritmus s nasledovným vstupom a výstupom (neformálne):

\begin{description}

\item[Vstup:] \begin{description}
\item referenčná postupnosť dusíkatých báz,
\item postupnosť nameraných hodnôt surového signálu,
\item \textit{nepovinné}: odhad očakávaného množstva variantov
\end{description}

\item[Výstup:] \begin{description}
\item popis nájdených variantov (pozícia, typ, skóre)
\end{description}
\end{description}

Náš algoritmus bude nevyhnutne robiť chyby. V niektorých prípadoch nezvládne nájsť variant, ktorý vzorka
obsahovala (falošné odmietnutie), v iných prípadoch nájde variant, ktorý neexistuje (falošné prijatie). 
V niektorých aplikáciách môže byť cena za chyby jedného druhu väčšia, než cena za chyby opačného 
druhu. Algoritmus preto vráti ku každému z nájdených variantov aj skóre, indikujúce istotu algoritmu, že 
naozaj ide o variant. Znižovaním minimálneho skóre, ktoré budeme vyžadovať, aby sme nájdený variant 
považovali za skutočný, bude možné znížit množstvo falošných odmietnutí za cenu zvýšenia množstva 
falošných prijatí, a obrátene.


Jedným z možných riešení nášho problému je určiť zo signálu bázy pomocou niektorého z existujúcich 
prekladačov báz a následne už len zisťovať odlišnosti dvoch postupností báz. Problémom tohoto riešenia
je nízka presnosť (ak nemáme veľa prekrývajúcich sa čítaní). 

Pri tomto prístupe však prekladač báz vôbec
nevyužíva fakt, že sekvenovaná postupnosť sa podobá na referenciu. Uvažuje teda podstatne väčší 
priestor možných výsledných sekvencií, než je nutné. V dôsledku toho upredňostňuje sekvencie, ktoré 
lepšie vysvetľujú pozorovaný signál, aj keď môžu byť výrazne vzdialené od referencie.
Zmenšenie priestoru uvažovaných sekvencií môže navyše znížiť výpočtovú náročnosť určovania báz, prípadne
umožniť použitie presnejších techník, ktoré by za normálnych okolností boli príliš časovo náročné.

V našej práci sa preto budeme snažiť nájsť riešenie, ktoré bude pri spustení na jednom čítaní presnejšie,
než porovnanie referencie s výstupom z prekladača báz.
